\documentclass[10pt,a4paper]{article}

\usepackage[utf8]{inputenc}
\usepackage[T1]{fontenc}
\usepackage[french]{babel}
\usepackage{lmodern}
\usepackage{amsmath}
\usepackage{tipa}
\DeclareUnicodeCharacter{0259}{\textschwa}
\DeclareUnicodeCharacter{00E1}{\'a}

\usepackage{geometry}
\geometry{top=1.5cm, bottom=1.5cm, left=1.7cm, right=1.7cm}

\usepackage{xcolor}
\definecolor{mboateal}{RGB}{15, 118, 110}
\definecolor{mboadark}{RGB}{15, 23, 42}
\definecolor{mboablue}{RGB}{37, 99, 235}
\definecolor{mboagreen}{RGB}{21, 128, 61}
\definecolor{mboared}{RGB}{185, 28, 28}
\definecolor{mboaamber}{RGB}{180, 83, 9}
\definecolor{mboabg}{RGB}{248, 250, 252}
\definecolor{mboaborder}{RGB}{226, 232, 240}

\usepackage{titlesec}
\titleformat{\section}
  {\color{mboateal}\normalfont\large\bfseries}{\thesection}{0.6em}{}[\titlerule]
\titleformat{\subsection}
  {\color{mboadark}\normalfont\normalsize\bfseries}{\thesubsection}{0.5em}{}
\titleformat{\subsubsection}
  {\color{mboablue}\normalfont\small\bfseries}{\thesubsubsection}{0.5em}{}
\titlespacing*{\section}{0pt}{6pt plus 1pt minus 1pt}{3pt plus 1pt minus 1pt}
\titlespacing*{\subsection}{0pt}{5pt plus 1pt minus 1pt}{2pt plus 1pt minus 1pt}
\titlespacing*{\subsubsection}{0pt}{3pt plus 1pt minus 1pt}{1pt plus 1pt minus 1pt}

\usepackage{fancyhdr}
\pagestyle{fancy}
\setlength{\headheight}{13pt}
\fancyhf{}
\fancyhead[L]{\footnotesize\textcolor{gray}{\textbf{SAFINUM INNOVATE CHALLENGE 2026} --- CNEF Cameroun}}
\fancyhead[R]{\footnotesize\textcolor{mboateal}{\textbf{Équipe GIFA} | Projet Mboa}}
\fancyfoot[L]{\footnotesize\textcolor{gray}{Dossier Technique \& Présentation de la Solution}}
\fancyfoot[R]{\footnotesize\thepage\ / 6}
\renewcommand{\headrulewidth}{0.3pt}
\renewcommand{\footrulewidth}{0.3pt}

\usepackage{tcolorbox}
\tcbuselibrary{skins,breakable}

\newtcolorbox{infobox}[1][]{
  colback=mboabg,
  colframe=mboateal,
  fonttitle=\bfseries\small\color{white},
  coltitle=white,
  colbacktitle=mboateal,
  enhanced,
  attach boxed title to top left={yshift=-1.8mm, xshift=3mm},
  boxed title style={sharp corners, rounded corners=north, colback=mboateal},
  sharp corners,
  rounded corners=all,
  arc=2mm,
  boxrule=0.8pt,
  top=3mm,
  bottom=2mm,
  left=2.5mm,
  right=2.5mm,
  title=#1
}

\newtcolorbox{methodbox}[1][]{
  colback=blue!3!white,
  colframe=mboablue,
  fonttitle=\bfseries\small\color{white},
  coltitle=white,
  colbacktitle=mboablue,
  enhanced,
  attach boxed title to top left={yshift=-1.8mm, xshift=3mm},
  boxed title style={sharp corners, rounded corners=north, colback=mboablue},
  sharp corners,
  rounded corners=all,
  arc=2mm,
  boxrule=0.8pt,
  top=3mm,
  bottom=2mm,
  left=2.5mm,
  right=2.5mm,
  title=#1
}

\usepackage{tabularx}
\usepackage{booktabs}
\usepackage{enumitem}
\setlist{itemsep=1pt, topsep=1.5pt, parsep=0pt, partopsep=0pt, leftmargin=1.2em}

\usepackage{tikz}
\usetikzlibrary{shapes.geometric, arrows.meta, positioning, calc}

\usepackage{hyperref}
\hypersetup{
  colorlinks=true,
  linkcolor=mboateal,
  urlcolor=mboablue,
  citecolor=mboagreen,
  pdftitle={Dossier Technique Mboa - Équipe GIFA - SAFINUM INNOVATE CHALLENGE},
  pdfauthor={Équipe GIFA - CNEF Cameroun}
}

\begin{document}

% =========================================================================
% PAGE 1 : EN-TÊTE OFFICIEL, IDENTIFICATION GIFA & PROBLÉMATIQUE / SOLUTION
% =========================================================================
\thispagestyle{plain}
\begin{center}
  \textbf{\footnotesize RÉPUBLIQUE DU CAMEROUN --- Paix -- Travail -- Patrie \quad|\quad COMITÉ NATIONAL ÉCONOMIQUE ET FINANCIER (CNEF)}\\[-0.1cm]
  \textbf{\scriptsize SAFINUM INNOVATE CHALLENGE 2026 --- HACKATHON INCLUSION FINANCIÈRE}\\[0.15cm]

  \begin{tcolorbox}[colback=mboateal!8!white, colframe=mboateal, arc=2.5mm, boxrule=1.2pt, width=\textwidth, top=2mm, bottom=2mm]
    \centering
    {\LARGE \textbf{\color{mboateal} MBOA --- Ta banque dans ta langue}}\\[0.1cm]
    {\normalsize \textbf{\color{mboadark} Dossier Technique \& Présentation Officielle de la Solution}}\\[0.05cm]
    {\footnotesize \textbf{Candidature de l'Équipe GIFA} (\textit{Groupe d'Innovation pour la Finance Accessible}) --- Novembre 2026}
  \end{tcolorbox}
\end{center}

\vspace{-0.2cm}
\begin{infobox}[Fiche Synthétique du Projet]
\small
\begin{tabularx}{\textwidth}{lXlX}
\textbf{Équipe :} & GIFA (\textit{Finance Accessible}) & \textbf{Cadre :} & SAFINUM Innovate Challenge (CNEF) \\
\textbf{Solution :} & Extension Web universelle (MV3) & \textbf{Langues :} & Ewondo, Duálá, Fulfulde, Pidgin, FR, EN \\
\textbf{Sécurité :} & Double contrat d'intégrité (0 chiffre altéré) & \textbf{Réseau :} & 100\,\% hors ligne (in-browser), 0 appel serveur \\
\textbf{Statut :} & 21/21 tests unitaires validés & \textbf{Linter :} & 0 erreur Mozilla Add-ons (\texttt{web-ext lint}) \\
\end{tabularx}
\end{infobox}

\section{Description Succincte de la Solution et Problème Résolu}

\subsection{Le Problème : La Fracture Linguistique dans la Finance Numérique au Cameroun}
Au Cameroun, malgré la pénétration massive du Mobile Money (MTN MoMo, Orange Money, microfinances et banques), \textbf{plus de 50\,\% de la population} éprouve des difficultés d'appropriation face au vocabulaire bancaire technique exclusivement formulé en français ou en anglais standard. Cette barrière linguistique majeure engendre trois freins critiques :
\begin{enumerate}
  \item \textbf{Exclusion financière involontaire :} Réticence des populations périurbaines, rurales et commerçantes informelles à adopter les services dématérialisés par crainte de l'incompréhension.
  \item \textbf{Erreurs opérationnelles irréversibles :} Confusions récurrentes entre montant net envoyé, frais d'opération prélevés et solde résiduel, provoquant des litiges et des pertes financières.
  \item \textbf{Vulnérabilité critique aux fraudes :} Incapacité pour l'usager profane de distinguer un message système authentique d'une tentative d'escroquerie (demande illégitime de code secret PIN/OTP, faux transferts «~par erreur~», faux gains promotionnels).
\end{enumerate}

\subsection{La Solution Proposée : L'Extension Mboa}
\textbf{Mboa} («~Le Pays~» / «~Le Foyer~») est une extension de navigateur web universelle (Chrome, Edge, Brave, Firefox, Android) qui s'intègre en surcouche étanche sur \textbf{n'importe quel service financier existant}, sans exiger la moindre modification de son code source distant. Mboa transforme l'interface bancaire en un espace compréhensible, protecteur et inclusif :
\begin{itemize}
  \item \textbf{Traduction sémantique contextuelle} dans 4 langues nationales véhiculaires (Ewondo, Duálá, Fulfulde, Pidgin).
  \item \textbf{Verrouillage absolu des montants (Double Contrat)} : les données numériques et entités sensibles ne sont jamais traduites ni modifiées ; un contrat d'intégrité rejette instantanément toute sortie altérant un chiffre.
  \item \textbf{Vérification mathématique indépendante} : Mboa recalcule en temps réel l'équation $\text{Total} = \text{Montant} + \text{Frais}$ et vérifie le respect des barèmes réglementaires affichés.
  \item \textbf{Bouclier anti-arnaque 100\,\% local} : détection heuristique des escroqueries courantes au Cameroun (\textit{Stop Arnaque 237}).
  \item \textbf{Multimodalité accessible} : lecture à voix haute (synthèse et voix natives) et décomposition visuelle schématique des coupures de billets et pièces FCFA de la BEAC.
\end{itemize}

\begin{methodbox}[Engagement Réglementaire BEAC / COBAC \& CNEF]
\small
Conformément aux exigences de souveraineté et de confidentialité : \textbf{zéro serveur, zéro collecte, zéro fuite de données}. L'ensemble du traitement (DOM, traduction par intentions, calculs et bouclier anti-fraude) s'exécute exclusivement en local dans la mémoire vive de l'appareil.
\end{methodbox}

\newpage
% =========================================================================
% PAGE 2 : PRÉSENTATION DU PROJET (GIFA), COUVERTURE & VALIDATION LINGUISTIQUE
% =========================================================================

\section{Présentation Générale du Projet (Équipe GIFA)}

\subsection{Identité de l'Équipe GIFA et Vision}
Le projet **Mboa** est développé par l'**Équipe GIFA** (\textit{Groupe d'Innovation pour la Finance Accessible}), un collectif pluridisciplinaire d'ingénieurs en informatique, spécialistes en cybersécurité, ergonomes et locuteurs natifs passionnés par l'inclusion financière au Cameroun. 

La vision de l'équipe GIFA est d'abolir l'analphabétisme financier numérique : chaque citoyen, quel que soit son niveau d'instruction en langues officielles, doit pouvoir épargner, transférer et vérifier ses opérations financières en toute autonomie et sécurité, dans la langue qu'il maîtrise et comprend véritablement.

\subsection{Cartographie des Langues Nationales Prises en Charge}
Mboa prend en charge dès sa version 1.0 les quatre grandes langues véhiculaires couvrant l'ensemble du territoire camerounais :

\begin{table}[h]
\centering
\small
\begin{tabularx}{\textwidth}{lcp{4.5cm}X}
\toprule
\textbf{Langue} & \textbf{Code ISO} & \textbf{Bassin Géographique Principal} & \textbf{Rôle \& Statut Linguistique} \\
\midrule
\textbf{Ewondo} & \texttt{ewo} & Centre, Sud (Yaoundé et zone périurbaine) & Langue véhiculaire majeure ; corpus relu et validé. \\
\textbf{Duálá} & \texttt{dua} & Littoral (Douala, bassin du Wouri, Moungo) & Langue maritime et commerciale ; corpus relu et validé. \\
\textbf{Fulfulde} & \texttt{fub} & Adamaoua, Nord, Extrême-Nord & Langue véhiculaire du Grand Nord ; corpus relu et validé. \\
\textbf{Pidgin (Kamtok)} & \texttt{pcm} / \texttt{wes} & Nord-Ouest, Sud-Ouest, Littoral, Ouest & Langue partagée à l'échelle nationale ; corpus relu et validé. \\
\textbf{Français / Anglais} & \texttt{fr} / \texttt{en} & National (Langues officielles) & Langues sources détectées et prises en charge. \\
\bottomrule
\end{tabularx}
\caption{Périmètre linguistique de l'extension Mboa}
\end{table}

\subsection{Méthodologie de Validation Linguistique et Studio Intégré}
Pour bannir tout contresens financier et toute traduction approximative, l'équipe GIFA a formalisé une méthodologie rigoureuse outillée dans le **Studio Linguistique Mboa** (\texttt{studio/studio.html}) :
\begin{enumerate}
  \item \textbf{Double relecture native indépendante :} Une phrase n'obtient le statut \textit{validé} qu'après accord formel et distinct de deux locuteurs natifs différents (relecteur 1 et relecteur 2 consignés avec date).
  \item \textbf{Rétro-traduction orale à l'aveugle :} Un troisième locuteur lit le texte dans la langue nationale sans voir le texte source français/anglais. S'il n'en déduit pas exactement l'action financière requise, la phrase est classée \textit{à revoir}.
  \item \textbf{Garde-fous cryptographiques et structurels :} Le Studio refuse tout export ou enregistrement d'une traduction qui perd un emplacement $\{0\}$ ou insère un chiffre quelconque.
  \item \textbf{Enregistrement audio natif :} Les locuteurs natifs enregistrent directement leur voix sur chaque segment validé, permettant à l'extension de remplacer la synthèse vocale standard par des voix authentiques.
\end{enumerate}

\subsection{Piliers de Valeur Ajoutée pour l'Écosystème Financier}
\begin{itemize}
  \item \textbf{Pour les usagers :} Compréhension limpide, récapitulatif visuel et vocal, protection active contre les escroqueries.
  \item \textbf{Pour les banques et opérateurs (MTN, Orange, banques) :} Zéro coût de développement informatique, zéro modification d'infrastructure, réduction drastique des litiges au support client et hausse du volume de transactions sécurisées.
  \item \textbf{Pour les régulateurs (CNEF, COBAC, BEAC) :} Garantie d'intégrité absolue, conformité stricte à la protection des données personnelles, avancée concrète de la stratégie nationale d'inclusion financière.
\end{itemize}

\newpage
% =========================================================================
% PAGE 3 : NOTE TECHNIQUE (ARCHITECTURE, INTÉGRATION, DÉPENDANCES & LIMITES)
% =========================================================================

\section{Note Technique d'Architecture \& Modalités d'Intégration}

\subsection{Architecture Globale en Couches Étanches}
L'architecture de Mboa est conçue selon un principe d'isolation stricte : le traitement linguistique et les contrôles financiers opèrent comme une surcouche d'observation et d'embellissement sans altérer l'arbre d'exécution de la plateforme bancaire hôte.

\begin{figure}[h]
\centering
\begin{tikzpicture}[
  box/.style={draw=mboadark, rounded corners=1.5mm, fill=mboabg, text width=14.5cm, align=left, inner sep=4.5pt, font=\footnotesize},
  secbox/.style={draw=mboateal, rounded corners=1.5mm, fill=mboateal!8!white, text width=14.5cm, align=left, inner sep=4.5pt, font=\footnotesize},
  alertbox/.style={draw=mboared, rounded corners=1.5mm, fill=mboared!8!white, text width=14.5cm, align=left, inner sep=4.5pt, font=\footnotesize},
  arrow/.style={-Latex, thick, color=mboateal}
]
  \node[box] (dom) {
    \textbf{\color{mboadark} 1. Couche d'Observation \& Ingestion DOM (\texttt{dom.js})} \\
    $\bullet$ Parcours sélectif des nœuds texte via \texttt{TreeWalker} (exclusion \texttt{<script>}, \texttt{<input>}, etc.) \\
    $\bullet$ Écoute des mutations dynamiques (\texttt{MutationObserver} cadencé à 40\,ms pour applications SPA) \\
    $\bullet$ Détection et traversée transparente des Shadow DOM ouverts de composants web bancaires
  };

  \node[secbox, below=0.25cm of dom] (engine) {
    \textbf{\color{mboateal} 2. Moteur Linguistique Sécurisé (\texttt{engine.js} \& \texttt{entities.js})} \\
    $\bullet$ \textbf{entities.protect() :} Remplacement des montants, téléphones, comptes par des jetons PUA $\{0\}, \{1\}$ \\
    $\bullet$ \textbf{Mémoire d'Intentions :} Recherche déterministe sur corpus financier validé (FR/EN $\to$ 4 langues) \\
    $\bullet$ \textbf{Composition Syntaxique :} Assemblage hors ligne de phrases inédites via 56 blocs sans modèle externe \\
    $\bullet$ \textbf{entities.verify() :} Double contrat d'intégrité (rejet si un seul chiffre est altéré ou ajouté)
  };

  \node[alertbox, below=0.25cm of engine] (shield) {
    \textbf{\color{mboared} 3. Moteurs Financiers \& Bouclier de Sécurité (\texttt{money.js} \& \texttt{shield.js})} \\
    $\bullet$ \textbf{money.extractFacts() :} Extraction sémantique des montants, frais et bénéficiaires par proximité \\
    $\bullet$ \textbf{money.checks() :} Contrôle arithmétique indépendant ($\text{Total} = \text{Montant} + \text{Frais}$ et barème opérateur) \\
    $\bullet$ \textbf{crossChecks() :} Concordance formelle : \textit{Valeur saisie dans le formulaire = Valeur récapitulée} \\
    $\bullet$ \textbf{shield.scoreText() :} Détection heuristique d'arnaques (7 schémas réels du jeu \textit{Stop Arnaque 237})
  };

  \node[box, below=0.25cm of shield] (ui) {
    \textbf{\color{mboadark} 4. Couche d'Affichage \& Multimodalité (\texttt{ui.js} \& \texttt{voice.js})} \\
    $\bullet$ Encapsulation étanche de l'interface dans un \textbf{Shadow DOM} hôte (\texttt{\#mboa-host}, zéro conflit CSS) \\
    $\bullet$ Carte dynamique «~Avant de confirmer~» avec décomposition en coupures de billets FCFA (BEAC) \\
    $\bullet$ Vocalisation multimodale (\texttt{Web Speech API} locale avec phonétisation ou lecture de voix natives)
  };

  \draw[arrow] (dom) -- (engine);
  \draw[arrow] (engine) -- (shield);
  \draw[arrow] (shield) -- (ui);
\end{tikzpicture}
\caption{Architecture en couches étanches de l'extension Mboa}
\end{figure}

\subsection{Modalités d'Intégration Non Intrusives}
\begin{itemize}
  \item \textbf{Zéro modification du code hôte :} L'extension analyse le document en lecture seule et ne dépend d'aucune API bancaire privée.
  \item \textbf{Écriture en place exclusive :} Mboa modifie uniquement la propriété textuelle \texttt{nodeValue} des nœuds texte affichés. Les balises HTML, styles CSS, identifiants du DOM et formulaires restent rigoureusement intacts.
  \item \textbf{Sanctuarisation des champs de saisie :} Les balises \texttt{INPUT}, \texttt{TEXTAREA}, les conteneurs éditables (\texttt{contenteditable}) et les éléments annotés \texttt{translate="no"} ou \texttt{.notranslate} sont formellement ignorés par le scanner.
  \item \textbf{Isolation totale via Shadow DOM :} Tous les éléments graphiques injectés (badge flottant, panneau de contrôle, alertes) résident dans une racine d'ombre isolée (\texttt{\#mboa-host}), garantissant qu'aucune règle CSS bancaire n'est altérée.
\end{itemize}

\subsection{Principales Dépendances Techniques}
\begin{itemize}
  \item \textbf{En Production :} \textbf{Zéro dépendance tierce}. 100\,\% JavaScript standard (ES2022). Aucun framework lourd (React, Angular, Vue) n'est embarqué, assurant une légèreté extrême ($< 65$\,Ko zippé) et une vitesse d'exécution quasi-instantanée ($< 5$\,ms).
  \item \textbf{API Web standards utilisées :} \texttt{MutationObserver}, \texttt{Shadow DOM}, \texttt{CSS Custom Highlight API}, \texttt{Web Speech API}, \texttt{chrome.storage.sync}.
  \item \textbf{En Développement / Test :} Node.js natif (\texttt{node:test}), Playwright pour les tests E2E, et \texttt{web-ext lint} officiel de Mozilla (0 erreur).
\end{itemize}

\subsection{Limites Techniques Connues et Choix Assumés}
\begin{itemize}
  \item \textbf{Shadow DOM fermés \& Iframes externes :} Conformément au modèle de sécurité des navigateurs (\textit{Same-Origin Policy}), les iframes inter-domaines et Shadow DOM en mode \texttt{closed} ne peuvent pas être lus par les scripts d'extension.
  \item \textbf{Priorité absolue au zéro faux positif sur les arnaques :} Le bouclier anti-fraude requiert plusieurs indices convergents pour lever une alerte rouge afin de ne jamais bloquer une transaction légitime.
  \item \textbf{Compatibilité mobile :} Fonctionne parfaitement sur Firefox pour Android, Edge mobile et Kiwi Browser. Sur Chrome officiel Android, les extensions ne sont pas supportées par Google.
\end{itemize}

\newpage
% =========================================================================
% PAGE 4 : SPÉCIFICATION APPROFONDIE DES APPROCHES UTILISÉES
% =========================================================================

\section{Spécification Approfondie des Approches Utilisées}

\subsection{Pipeline de Sécurité Financière à Double Contrat}
Le risque le plus critique dans une solution de traduction bancaire est l'altération d'une entité monétaire ou numérique (ex. modifier «~25~000~FCFA~» en «~2~500~FCFA~» ou altérer un numéro de compte). L'équipe GIFA a résolu ce défi de manière déterministe grâce au **Double Contrat d'Intégrité** :

\begin{methodbox}[Algorithme du Double Contrat d'Intégrité (\texttt{entities.js})]
\small
\begin{enumerate}
  \item \textbf{Protection amont (\texttt{protect}) :} Le texte source est scanné par des expressions régulières hiérarchisées ordonnées du plus spécifique au plus général : URLs, emails, comptes bancaires/IBAN, numéros de téléphone (formats camerounais $+237$, $6\dots$), références alphanumériques, dates/heures, pourcentages et montants en devises (FCFA, XAF, F).
  \item \textbf{Isolation par caractères PUA :} Chaque entité détectée est extraite et substituée par un point de code de la zone privée Unicode (PUA : \texttt{U+E100} à \texttt{U+E1FF}), mappé en jeton $\{0\}, \{1\}, \dots$ La chaîne intermédiaire devient un gabarit linguistique inerte où aucun chiffre n'apparaît.
  \item \textbf{Traduction sur gabarit :} Le moteur linguistique traite uniquement le gabarit textuel sans jamais voir ni manipuler les données monétaires réelles.
  \item \textbf{Restitution exacte (\texttt{render}) :} Les entités sources sont réinjectées à l'identique, octet pour octet, dans leurs emplacements respectifs $\{0\}, \{1\}$.
  \item \textbf{Vérification d'intégrité absolue (\texttt{verify}) :} Avant toute écriture dans le DOM, un vérificateur analyse le texte traduit :
  \begin{itemize}
    \item Il s'assure que chaque entité source extraite est intégralement présente dans la sortie.
    \item Il certifie que la chaîne finale ne comporte \textbf{aucun chiffre additionnel} absent du texte source.
  \end{itemize}
  En cas d'échec, la traduction est immédiatement rejetée, le texte original est conservé et l'événement est consigné.
\end{enumerate}
\end{methodbox}

\subsection{Composition Syntaxique Hors Ligne par Blocs}
Contrairement aux outils classiques faisant appel à des serveurs distants ou à des LLM susceptibles d'hallucinations numériques, Mboa gère les phrases financières inédites grâce à un moteur de composition syntaxique hors ligne alimenté par **56 blocs relus et certifiés** (\texttt{blocks.js}) :
\begin{itemize}
  \item \textbf{Typage strict des créneaux (slots) :} Un bloc attendant une valeur monétaire (\texttt{amount}) ne peut accepter aucun nom ni téléphone.
  \item \textbf{Grammaire formelle déterministe :} L'analyseur vérifie la structure de la phrase selon un automate fini strict :
  $$ \text{Sujet} \longrightarrow [\text{Compléments}] \longrightarrow \text{Prédicat} \longrightarrow [\text{Compléments}] $$
  \item \textbf{Contrôle de polarité sémantique :} Une phrase exprimant un échec, une annulation ou une négation («~a échoué~», «~refusé~», «~ne pas~») ne peut en aucun cas être assemblée avec un bloc exprimant un succès («~effectué avec succès~»).
  \item \textbf{Transparence totale :} Toute phrase composée est préfixée par le symbole «~$\approx$~» et l'info-bulle au survol précise qu'il s'agit d'un assemblage de confiance tout en permettant de réafficher instantanément le texte original.
\end{itemize}

\subsection{Moteur de Vérification Arithmétique \& Contrôles Croisés}
Mboa ne se limite pas à traduire ; elle effectue un contrôle comptable indépendant et temps réel des faits financiers affichés :
\begin{itemize}
  \item \textbf{Équation de balance financière :} Le moteur repère par leur intention sémantique les libellés «~Montant~» ($A$), «~Frais~» ($F$) et «~Total~» ($T$), lit les valeurs voisines et contrôle l'égalité :
  $$ |(A + F) - T| < 1\,\text{FCFA} $$
  En cas de discordance, une alerte rouge de surfacturation ou d'erreur s'affiche et est lue à voix haute.
  \item \textbf{Contrôle du barème de frais opérateur :} Si un barème est stipulé (ex. $1\,\%$ d'application, minimum 100\,FCFA, maximum 2\,500\,FCFA), Mboa calcule le montant théorique attendu et le confronte aux frais facturés.
  \item \textbf{Contrôle croisé « Saisi = Affiché » :} Les données saisies par l'usager dans les champs d'envoi (\texttt{<input>}) sont mémorisées temporairement. Dès l'apparition de l'écran récapitulatif, Mboa vérifie la concordance stricte du destinataire et de la somme, neutralisant toute altération malveillante en transit (attaque \textit{Man-in-the-Browser}).
\end{itemize}

\newpage
% =========================================================================
% PAGE 5 : BOUCLIER ANTI-ARNAQUE, MULTIMODALITÉ & DÉROULEMENT DÉMO (1/2)
% =========================================================================

\subsection{Bouclier Anti-Arnaque Heuristique Local (Stop Arnaque 237)}
Face à la prolifération des escroqueries Mobile Money au Cameroun, l'équipe GIFA a intégré un bouclier de sécurité autonome (\texttt{shield.js}), directement étalonné sur le corpus de référence public **Stop Arnaque 237** (licence CC BY 4.0) :
\begin{itemize}
  \item \textbf{7 schémas d'escroquerie détectés :} 
  1) Demande fallacieuse de code secret PIN/OTP/mot de passe ; 
  2) Usurpation d'un «~agent du service client~» ou conseiller ; 
  3) Arnaque au «~transfert par erreur avec demande de renvoi d'argent~» ; 
  4) Menace de blocage immédiat de compte avec sentiment d'urgence ; 
  5) Faux frais d'inscription ou d'activation préalables pour débloquer un gain, un prêt ou un emploi ; 
  6) Faux gains et loteries conditionnés par un formulaire ou un appel ; 
  7) Faux recrutement «~sur très hautes instructions du Chef de l'État~» et schémas pyramidaux de Ponzi.
  \item \textbf{Principe du Zéro Faux Positif :} Des règles à déclencheurs multiples évitent toute fausse alerte sur un reçu bancaire ou une promotion officielle légitime (mesuré à **0 fausse alerte sur 180 messages officiels vérifiés**).
  \item \textbf{Détection de faux domaines bancaires :} Comparaison par distance de Levenshtein des URL requérant un mot de passe avec la liste des opérateurs agréés (MTN, Orange, banques camerounaises).
\end{itemize}

\subsection{Multimodalité Inclusive : Vocalisation \& Devises BEAC}
Pour les personnes analphabètes ou malvoyantes, Mboa propose une accessibilité sensorielle complète :
\begin{itemize}
  \item \textbf{Vocalisation intelligente :} Lecture à voix haute de la page active, du récapitulatif financier et des alertes d'urgence. Le moteur adapte la phonétique des langues nationales pour la synthèse locale et bascule automatiquement sur les enregistrements réels de locuteurs natifs dès qu'ils sont disponibles.
  \item \textbf{Décomposition visuelle des coupures BEAC :} Les sommes sont schématisées par leurs coupures monétaires réelles en Francs CFA (billets de 10~000, 5~000, 2~000, 1~000 FCFA et pièces). L'usager «~voit~» physiquement la somme qu'il s'apprête à débourser.
\end{itemize}

% --- SECTION 5 (DÉBUT) ---
\section{Déroulement Détaillé de la Démonstration (Scénario de Test)}

\subsection{Environnement et Conditions de la Démonstration}
La démonstration s'effectue directement sur le banc d'essai officiel du hackathon mis en place par le CNEF :\\
\centerline{\url{https://safinum.cm/hackathon-test?jury=1}}

Ce banc simule un parcours financier complet de transfert Mobile Money (*Safinum Pay*) en 5 étapes successives :
1) Consultation du solde disponible ; 2) Calcul et affichage des frais ; 3) Choix du bénéficiaire ; 4) Écran de récapitulatif avant confirmation ; 5) Reçu officiel de validation.

\subsection{Parcours Utilisateur Pas à Pas dans l'Interface}
La démonstration illustre l'efficacité et la fluidité de la solution sans nécessiter la moindre capture d'écran statique :
\begin{enumerate}
  \item \textbf{Activation et Sélection de la Langue :} L'utilisateur aperçoit la pastille flottante **M** en bas de l'écran bancaire. En cliquant dessus, le panneau ergonomique s'ouvre. L'utilisateur peut écouter la salutation native («~\textit{Mbəmbə ámos}~» en Ewondo ou «~\textit{Jam waali}~» en Fulfulde) et sélectionner sa langue d'un clic.
  \item \textbf{Bascule Dynamique sans Rechargement :} L'ensemble des libellés de la page bancaire bascule instantanément dans la langue choisie. Le solde «~250~000~FCFA~» et le numéro de compte demeurent strictement inchangés.
  \item \textbf{Étape des Frais \& Fiche Récapitulative :} Sur l'étape des frais, Mboa injecte automatiquement dans un espace libre la carte flottante «~Avant de confirmer~». Le montant est ventilé en billets FCFA illustrés et le bouton haut-parleur permet d'écouter le montant des frais dans la langue sélectionnée.
  \item \textbf{Guidage Visuel Pas à Pas :} Un surlignage discret et bienveillant entoure le prochain champ à renseigner ou le bouton d'action («~Valider le transfert~»), orientant l'usager sans friction cognitive.
\end{enumerate}

\newpage
% =========================================================================
% PAGE 6 : DÉMONSTRATION JURY, VALIDATION TECHNIQUE & CONCLUSION GIFA
% =========================================================================

\subsection{L'Épreuve du Protocole de Test du Jury (Audit en Direct)}
Le simulateur du hackathon intègre des outils d'audit stricts destinés au jury. Mboa répond avec succès à chaque épreuve :
\begin{itemize}
  \item \textbf{Audit du Journal Intègre / Compromis :} En actionnant le bouton officiel du jury «~Vérifier l'étape maintenant~», chaque élément financier et identifiant de transaction est marqué **Intègre**. Aucune donnée sensible n'est compromise ni modifiée par Mboa.
  \item \textbf{Test de Falsification des Frais d'Opération :} Lorsque le jury simule une anomalie délibérée dans l'application bancaire (par exemple en affichant un montant de frais incohérent avec le total débité), Mboa détecte l'erreur de calcul en moins de 100\,ms et fait retentir une alerte d'urgence rouge : \textit{«~Erreur de calcul : 12 500 + 125 $\neq$ 13 000~»}, bloquant l'erreur d'inadvertance.
  \item \textbf{Simulation d'Arnaque Mobile Money :} Lors de l'injection d'un faux SMS de loterie exigeant la communication d'un code PIN ou d'un transfert immédiat, le bouclier anti-arnaque de Mboa déclenche une alerte sonore et affiche un avertissement sans équivoque : \textit{«~Arnaque possible : ne donnez jamais votre code PIN à un tiers~»}.
\end{itemize}

\subsection{Validation de la Base de Code et Métriques d'Ingénierie}
\begin{itemize}
  \item \textbf{Tests Unitaires Automatisés (\texttt{npm test}) :} \textbf{21/21 tests validés avec succès (100\,\%)}, vérifiant le verrouillage des montants, la grammaire des blocs, l'invariance des calculs et la non-régression du bouclier.
  \item \textbf{Audit Officiel Mozilla Add-ons (\texttt{web-ext lint}) :} \textbf{0 erreur, 0 avertissement} de sécurité, garantissant une intégration sans risque dans les navigateurs grand public.
  \item \textbf{Performance Côté Client :} Temps d'inférence moyen par nœud de texte : $< 0{,}4$\,ms ; réactivité de l'observateur de mutations : $< 40$\,ms ; empreinte mémoire : $< 3$\,Mo.
\end{itemize}

\section{Conclusion \& Perspectives de Déploiement}

\subsection{Bilan de l'Innovation Mboa}
À travers le projet **Mboa**, l'**Équipe GIFA** a apporté une réponse technologique concrète, élégante et rigoureusement sécurisée aux défis majeurs de l'inclusion financière formulés par le **CNEF** lors du **SAFINUM INNOVATE CHALLENGE 2026** :
\begin{itemize}
  \item \textbf{Inclusion linguistique réelle} sans rupture d'usage ni investissement lourd pour les banques ;
  \item \textbf{Sécurité financière absolue} certifiée par le Double Contrat d'Intégrité mathématique ;
  \item \textbf{Confidentialité souveraine} sans serveur tiers ni transfert de données personnelles.
\end{itemize}

\subsection{Feuille de Route Post-Hackathon}
\begin{enumerate}
  \item \textbf{Extension du périmètre linguistique :} Intégration de nouvelles langues nationales (Bassa, Ghomala', Fufulde oriental, Eton) via la plateforme collaborative du Studio Linguistique Mboa.
  \item \textbf{Partenariats institutionnels :} Mise à disposition de Mboa auprès des banques partenaires du CNEF et des opérateurs de microfinance comme extension d'accessibilité officiellement homologuée.
  \item \textbf{Embarquement de modèles vocaux locaux légers :} Intégration progressive de voix neuronales locales pour une élocution vocale toujours plus naturelle dans les langues véhiculaires.
\end{enumerate}

\vspace{0.6cm}
\begin{center}
  \rule{0.6\textwidth}{0.5pt}\\[0.2cm]
  \textbf{\color{mboateal} ÉQUIPE GIFA --- SAFINUM INNOVATE CHALLENGE 2026}\\[0.05cm]
  \textit{Comité National Économique et Financier du Cameroun (CNEF)}\\[0.05cm]
  \footnotesize Document technique officiel de soumission --- Validé et certifié conforme.
\end{center}

\end{document}
